% Chapter 4, Section 2 _Linear Algebra_ Jim Hefferon
%  http://joshua.smcvt.edu/linearalgebra
%  2001-Jun-12
\section{行列式的几何}
前面一节从代数角度展开行列式的讨论，考虑的是满足若干条件的公式。
本节则以几何视角作为补充。
这种做法除了直观上的吸引力之外，还有一个优点：到目前为止我们只关心
行列式是否为零，而在这里我们将赋予行列式的值以含义。
（前面一节把行列式视为行的函数，而本节聚焦于列。）










\subsection{作为大小函数的行列式}
这幅平行四边形的图形，在构造两个向量之和时已为大家所熟悉。
\begin{center}
  \includegraphics{det/mp/ch4.30}
\end{center}

\begin{definition} \label{df:Box}
%<*df:Box>
在 $\Re^n$ 中，
由 \( \sequence{\vec{v}_1,\dots,\vec{v}_n} \) 生成的
\definend{箱体}\index{box}
（或 \definend{平行六面体}\index{parallelepiped}）
是集合
\( \set{t_1\vec{v}_1+\dots+t_n\vec{v}_n
      \suchthat t_1,\ldots,t_n\in \closedinterval{0}{1}} \)。
%</df:Box>
\end{definition}

\noindent 于是，上面的平行四边形就是
由 $\sequence{\binom{x_1}{y_1},\binom{x_2}{y_2}}$ 生成的箱体。
三维空间中的箱体见 \nearbyexample{ex:VolParPiped}。

我们可以通过画一个包围矩形，再把不属于箱体的各块面积减去，
来求出上述箱体的面积。
\begin{center}
  \parbox{1.5in}{\hbox{}\hfil\includegraphics{det/mp/ch4.31}\hfil\hbox{}}
  \quad
  \parbox{3.0in}{
    \hbox{}\hfil
    $\begin{array}{l}
        \text{平行四边形的面积}                    \\
          \hbox{}\quad \hbox{}
             =\text{矩形的面积}
                 -\text{$A$ 的面积}-\text{$B$ 的面积} \\
          \hbox{}\qquad \hbox{}
             -\cdots-\text{$F$ 的面积}                   \\
          \hbox{}\quad \hbox{}
             =(x_1+x_2)(y_1+y_2)-x_2y_1-x_1y_1/2        \\
          \hbox{}\qquad \hbox{}
             -x_2y_2/2-x_2y_2/2-x_1y_1/2-x_2y_1         \\
          \hbox{}\quad \hbox{}
             =x_1y_2-x_2y_1
    \end{array}$
    \hfil\hbox{}}
\end{center}
该面积等于行列式
\begin{equation*}
  \begin{vmat}
    x_1  &x_2  \\
    y_1  &y_2
  \end{vmat}
  =x_1y_2-x_2y_1
\end{equation*}
的值，这并非巧合。
行列式的定义中包含四条性质，我们知道对于每个维数~$n$，
这些性质唯一确定一个函数。
我们将论证，这些性质恰好是一个度量 $n$ 维空间中箱体大小的函数
的合理公设。\index{size}

举例来说，这样的函数应当具有如下性质：把定义箱体的某个向量乘以
一个标量，
箱体的大小也随之乘以该标量。
\begin{center}
  \includegraphics{det/mp/ch4.32}
  \qquad
  \includegraphics{det/mp/ch4.33}
\end{center}
图中展示的是 $k=1.4$ 的情形。
右图中，缩放后的区域用实线画出，原来的区域涂以阴影以便比较。
% The region formed by $k\vec{v}$ and~$\vec{w}$
% is bigger by a factor of \( k \)
% than the shaded region enclosed by $\vec{v}$ and~$\vec{w}$.
% That is,
% \( \size (k\vec{v},\vec{w})=k\cdot\size (\vec{v},\vec{w}) \) and

也就是说，我们可以合理地预期
$\size (\dots,k\vec{v},\dots)=k\cdot\size (\dots,\vec{v},\dots)$。
当然，这一条件正是行列式定义中的性质之一。

行列式的另一条性质也应适用于任何度量箱体大小的函数：
它不受倍加操作的影响。
这里是倍加之前与之后的箱体（图中标出的标量为 $k=-0.35$）。
\begin{center}
  \includegraphics{det/mp/ch4.34}
  \qquad
  \includegraphics{det/mp/ch4.35}
\end{center}
由 $v$ 与
$k\vec{v}+\vec{w}$ 生成的箱体倾斜方式
与原来的不同，但二者具有相同的底和相同的高，因而面积相同。
% (As before, the figure on the right has the
% original region in shade for comparison.)
所以我们预期，大小不受剪切操作的影响：
$\size (\dots,\vec{v},\dots,\vec{w},\dots)
=\size (\dots,\vec{v},\dots,k\vec{v}+\vec{w},\dots)$。
同样，这也是一条行列式性质。

我们预期由单位向量生成的箱体具有单位大小
\begin{center}
  \includegraphics{det/mp/ch4.36}
\end{center}
并且自然地把它推广到任意 $n$ 维空间：
$\size(\vec{e}_1,\dots,\vec{e}_n)=1$。
% Again, that is a determinant property.

正如定义之后所指出的，行列式定义中的条件~(2) 是冗余的。
从前一节我们知道，对每个~$n$ 行列式都存在且唯一，
因此这些关于大小函数的公设是相容的，
而且我们不需要更多的公设。
于是，我们把 \( \det(\vec{v}_1,\dots,\vec{v}_n) \)
解释为给出生成箱体的向量所构成箱体的大小，就是有根据的。
% (\textit{Comment.}
%   An even more basic approach, which also leads to the definition
%   below, is in \cite{Weston59}.)

\begin{remark}  \label{re:PropertyTwoGivesSign}
虽然条件~(2) 是冗余的，但它提出了一个重要问题。
考虑下面两个行列式。
\begin{center} \small
  \begin{tabular}{c@{\hspace*{8em}}c}
    \includegraphics{det/mp/ch4.37}
      &\includegraphics{det/mp/ch4.38}  \\[.25ex]
    \ $\begin{vmat}[r]
        4  &1   \\
        2  &3
      \end{vmat}=10$
      &\ $\begin{vmat}[r]
          1  &4   \\
          3  &2
        \end{vmat}=-10$
  \end{tabular}
\end{center}
对换两列会改变符号。
在左图中，从 $\vec{u}$ 出发沿夹角内部的弧走到 $\vec{v}$（即沿逆时针方向），
我们得到正的大小。
在右图中，从 $\vec{v}$ 出发走到~$\vec{u}$，
即沿顺时针方向的弧，则得到负的大小。
大小函数所给出的符号反映了箱体的
\definend{定向}\index{box!orientation}\index{orientation}
或 \definend{朝向}\index{box!sense}\index{sense}。
（如果我们设想用一个负标量做标量乘法的效果，也会看到同样的现象。）
% Although it is both interesting and important, we don't need the idea of
% orientation for the development below and so we will pass it by.
% (See \nearbyexercise{exer:BasisOrient}.)
\end{remark}

\begin{definition} \label{df:Volume}
%<*df:Volume>
箱体的\definend{体积}\index{volume}\index{box!volume}
是以这些向量为列的矩阵的行列式的绝对值。
%</df:Volume>
\end{definition}

\begin{example} \label{ex:VolParPiped}
按底面积乘以高的公式，这个
平行六面体的体积是 $12$。
这与行列式的值一致。
\begin{center}
   \parbox{2in}{\hbox{}\hfil\includegraphics{det/asy/ppiped.pdf}\hfil\hbox{}}
  %  \parbox{2in}{\hbox{}\hfil\includegraphics{det/mp/ch4.39}\hfil\hbox{}}
  \hspace{4.5em}
  $\begin{vmat}[r]
     2 &0 &-1\\
     0 &3 &0 \\
     2 &1 &1
  \end{vmat}=12$
\end{center}
按不同次序取这些向量会改变符号，但不改变大小。
\begin{equation*}
  \begin{vmat}[r]
     0  &2 &-1 \\
     3  &0 &0 \\
     1  &2 &1
  \end{vmat}=-12
\end{equation*}
\end{example}

% The next result describes some of the geometry of the linear
% functions that act on
% \( \Re^n \).

\begin{theorem}\label{th:MatChVolByDetMat}
\index{size}\index{transformation!size change}
%<*th:MatChVolByDetMat>
变换 \( \map{t}{\Re^n}{\Re^n} \) 以同一个因子改变所有箱体的大小，也就是说，箱体的像的大小
$\deter{t(S)}$ 等于 $\deter{T}$ 乘以箱体的大小 $\deter{S}$，
其中 $T$ 是在标准基下
表示 $t$ 的矩阵。

也就是说，乘积的行列式等于
行列式的乘积 $\deter{TS}=\deter{T}\cdot\deter{S}$。
%</th:MatChVolByDetMat>
\end{theorem}

这两句话说的是同一件事，前者用映射的语言，后者用
矩阵的语言。
这是因为
$\deter{t(S)}=\deter{TS}$，二者给出的都是单位箱体 $\stdbasis_n$
在复合映射 $\composed{t}{s}$ 下的像这一箱体的大小，
其中各映射均以标准基下的矩阵表示。
我们将证明第二句话。

\begin{proof}
%<*pf:MatChVolByDetMat0>
首先考虑 $T$ 是奇异的、因而
没有逆矩阵的情形。
注意到，若 \( TS \) 可逆，则存在
某个 $M$ 使得 \( (TS)M=I \)，于是
\( T(SM)=I \)，从而 \( T \) 可逆。
这一论断的逆否命题是：若 \( T \) 不可逆，则 \( TS \) 也不可逆 \Dash
若 $\deter{T}=0$ 则 $\deter{TS}=0$。
%</pf:MatChVolByDetMat0>

%<*pf:MatChVolByDetMat1>
再考虑 $T$ 是非奇异的情形。
任何非奇异矩阵都可以分解为
初等矩阵的乘积 $T=E_1E_2\cdots E_r$。
为完成这一论证，
我们将对所有的矩阵~$S$ 与初等矩阵~$E$ 验证
\( \deter{ES}=\deter{E}\cdot\deter{S} \)。
于是结论随之成立，因为
$\deter{TS}=\deter{E_1\cdots E_rS}=\deter{E_1}\cdots\deter{E_r}\cdot\deter{S}
  =\deter{E_1\cdots E_r}\cdot\deter{S}=\deter{T}\cdot\deter{S}$。
%</pf:MatChVolByDetMat1>

%<*pf:MatChVolByDetMat2>
初等矩阵有三种类型。
我们处理 $M_i(k)$ 的情形；
$P_{i,j}$ 与 $C_{i,j}(k)$ 的验证类似。
矩阵 $M_i(k)S$ 等于 $S$，只是第~$i$ 行乘了 $k$。
由行列式函数的第三条条件
即得 $\deter{M_i(k)S}=k\cdot\deter{S}$。
而 $\deter{M_i(k)}=k$，这同样由第三条条件得出，因为
$M_i(k)$ 是把单位矩阵的第~$i$ 行乘以
$k$ 得到的。
因此 \( \deter{ES}=\deter{E}\cdot\deter{S} \) 对
$E=M_i(k)$ 成立。
%</pf:MatChVolByDetMat2>
\end{proof}


\begin{example}
应用这个以标准基下的矩阵
\begin{equation*}
  \begin{mat}[r]
    1  &1  \\
   -2  &0
  \end{mat}
\end{equation*}
表示的映射 $t$，会使箱体的大小加倍，例如，从这个
\begin{center}
    \parbox{1.5in}{\hbox{}\hfil\includegraphics{det/mp/ch4.40}\hfil\hbox{}}
    \quad
    $\begin{vmat}[r]
      2  &1  \\
      1  &2
    \end{vmat}=3$
\end{center}
到这个
\begin{center}
    \parbox{1.5in}{\hbox{}\hfil\includegraphics{det/mp/ch4.41}\hfil\hbox{}}
    \quad
    $\begin{vmat}[r]
        3  &3  \\
       -4  &-2
     \end{vmat}=6$
\end{center}
\end{example}


% Recall that determinants are not additive homomorphisms, that
% $\det(A+B)$ need not equal $\det(A)+\det(B)$.
% In contrast, the above theorem says that determinants are
% multiplicative homomorphisms:
% $\det(AB)$ equals $\det(A)\cdot \det(B)$.

\begin{corollary} \label{co:DeterminantOfInverseIsInverseOfDeterminant}
%<*co:DeterminantOfInverseIsInverseOfDeterminant>
若矩阵可逆，则其逆矩阵的行列式是其行列式的倒数
$\deter{T^{-1}}=1/\deter{T}$。
%</co:DeterminantOfInverseIsInverseOfDeterminant>
\end{corollary}

\begin{proof}
%<*pf:DeterminantOfInverseIsInverseOfDeterminant>
$1=\deter{I}=\deter{TT^{-1}}=\deter{T}\cdot\deter{T^{-1}}$
%</pf:DeterminantOfInverseIsInverseOfDeterminant>
\end{proof}


\begin{exercises}
  \item 
    问：\begin{equation*}
      \colvec[r]{4 \\ 1 \\ 2}
    \end{equation*}
    是否位于由下面这三个向量所张成的箱体内部？
    \begin{equation*}
      \colvec[r]{3 \\ 3 \\ 1}
      \quad
      \colvec[r]{2 \\ 6 \\ 1}
      \quad
      \colvec[r]{1 \\ 0 \\ 5}
    \end{equation*}
    \begin{answer}
      求解
      \begin{equation*}
        c_1\colvec[r]{3 \\ 3 \\ 1}
        +c_2\colvec[r]{2 \\ 6 \\ 1}
        +c_3\colvec[r]{1 \\ 0 \\ 5}
        =\colvec[r]{4 \\ 1 \\ 2}
      \end{equation*}
      得到唯一解
      \( c_3=11/57 \)、\( c_2=-40/57 \) 以及 \( c_1=99/57 \)。
      由于 \( c_1>1 \)，该向量不在箱体内。  
    \end{answer}
  \recommended \item 
    求由下列向量所定义的区域的体积。
    \begin{exparts}
      \partsitem $\sequence{\colvec[r]{1 \\ 3},\colvec[r]{-1 \\ 4}}$
      \partsitem $\sequence{\colvec[r]{2 \\ 1 \\ 0},\colvec[r]{3 \\ -2 \\ 4},
                              \colvec[r]{8 \\ -3 \\ 8}}$
      \partsitem $\sequence{\colvec[r]{1 \\ 2 \\ 0 \\ 1},
                             \colvec[r]{2 \\ 2 \\ 2 \\ 2},
                             \colvec[r]{-1 \\ 3 \\ 0 \\ 5},
                             \colvec[r]{0 \\ 1 \\ 0 \\ 7}}$
    \end{exparts}
    \begin{answer}
      对每一小题，求出行列式并取绝对值。
      \begin{exparts*}
        \partsitem $7$
        \partsitem $0$
         % sage: A = matrix(QQ, [[2, 3, 8], [1, -2, -3], [0, 4, 8]])
         % sage: A
         % [ 2  3  8]
         % [ 1 -2 -3]
         % [ 0  4  8]
         % sage: det(A)
         % 0
        \partsitem $-58$
         % sage: A = matrix(QQ, [[1,2,-1,0], [2,2,3,1], [0,2,0,0], [1,2,5,7]])
         % sage: A
         % [ 1  2 -1  0]
         % [ 2  2  3  1]
         % [ 0  2  0  0]
         % [ 1  2  5  7]
         % sage: det(A)
         % -58
      \end{exparts*}
    \end{answer}
  \recommended \item 
    在下图中，左边的矩形由点 $(3,1)$ 和
    $(1,2)$ 确定。  
    用该矩阵作用后得到右边的矩形，它由
    $(7,1)$ 和 $(4,2)$ 确定。
    为什么这并不与
    \nearbytheorem{th:MatChVolByDetMat} 矛盾？
    \begin{center}
      \begin{tabular}{c@{\hspace*{2em}}c@{\hspace*{2em}}c}
        \includegraphics{det/mp/ch4.43}
        &\raisebox{12pt}{\( \grstep{\bigl(\begin{smallmatrix}
                                        2  &1 \\
                                        0  &1
                                     \end{smallmatrix}\bigr)} \)}
        &\includegraphics{det/mp/ch4.44}                                 \\
        面积为 $2$
        &行列式为 $2$
        &面积为 $3$
      \end{tabular}
    \end{center}
    \begin{answer}
      我们画这幅图就是为了误导读者。
      左图并不是由两个向量形成的箱体。
      若把它平移到原点，它就变成由下面这个序列形成的箱体。
      \begin{equation*}
        \sequence{
          \colvec[r]{0 \\ 1},
          \colvec[r]{2 \\ 0}
       }
      \end{equation*}
      而矩阵作用后的像是由下面这个序列形成的箱体。
      \begin{equation*}
        \sequence{
          \colvec[r]{1 \\ 1},
          \colvec[r]{4 \\ 0}
         }
      \end{equation*}
      其面积为 $4$。
     \end{answer}
  \recommended \item 
    求下述区域的体积。
    \begin{center}
      \includegraphics{det/mp/ch4.42}  
    \end{center}
    \begin{answer}
      将该平行四边形平移至从原点出发，
      它就变成由下列向量形成的箱体
      \begin{equation*}
        \sequence{
          \colvec[r]{3 \\ 0},
          \colvec[r]{2 \\ 1}
        }      
      \end{equation*}
      于是这个行列式的绝对值
      容易算出为 $3$。
      \begin{equation*}
        \begin{vmat}[r]
          3  &2  \\
          0  &1
        \end{vmat}=3
      \end{equation*}
     \end{answer}
  \recommended \item 
    设 \( \deter{A}=3 \)。
    下列各量把体积改变了多少倍？
    \begin{exparts*}
      \partsitem \( A \)
      \partsitem \( A^2 \)
      \partsitem \( A^{-2} \)
    \end{exparts*}
    \begin{answer}
     \begin{exparts*}
        \partsitem \( 3 \)
        \partsitem \( 9 \)
        \partsitem $1/9$
      \end{exparts*}  
    \end{answer}
  \recommended \item
    考虑~$\Re^3$ 上的线性变换~$t$，它关于
    标准基由下述矩阵表示。
    \begin{equation*}
      \begin{mat}
        1 &0 &-1 \\
        3 &1 &1 \\
       -1 &0 &3
      \end{mat}
    \end{equation*}
    \begin{exparts}
      \partsitem 计算该矩阵的行列式。
        这个变换是保持定向还是反转定向？ 
      \partsitem 求由这些向量所定义的箱体的大小。
        它的定向是什么？
        \begin{equation*}
          \colvec{1 \\ -1 \\ 2}
          \quad
          \colvec{2 \\ 0 \\ -1}
          \quad
          \colvec{1 \\ 1 \\ 0}
        \end{equation*}
      \partsitem 求前一小题中各向量在 $t$ 下的像，
        并求出它们所定义的箱体的大小。
        定向是什么？
    \end{exparts}
    \begin{answer}
      \begin{exparts}
        \partsitem
          用高斯消元法
          \begin{equation*}
          \grstep[\rho_1+\rho_3]{-3\rho_1+\rho_2}
          \begin{mat}
            1 &0   &-1    \\
            0 &1   &4     \\
            0 &0   &2
          \end{mat}
          \end{equation*}
          得到行列式为~$+2$。
          符号为正，所以该变换保持定向。
        \partsitem
          箱体的大小就是下述行列式的值。
          \begin{equation*}
            \begin{vmat}
              1 &2  &1 \\
             -1 &0  &1 \\
              2 &-1 &0
            \end{vmat}
            =+6
          \end{equation*}
          定向为正。
        \partsitem
          由于该变换关于标准基由所给矩阵表示，
          而关于标准基，向量就表示它们自身，
          所以要求这些向量在该变换下的像，
          只需从左边用矩阵去乘它们。
          \begin{equation*}
            \colvec{1 \\ -1 \\ 2}\mapsto\colvec{-1 \\ 4 \\ 5}
            \qquad
            \colvec{2 \\ 0 \\ -1}\mapsto\colvec{3 \\ 5 \\ -5}
            \qquad
            \colvec{1 \\ 1 \\ 0}\mapsto\colvec{1 \\ 4 \\ -1}
          \end{equation*}
          然后计算所得箱体的大小。
          \begin{equation*}
            \begin{vmat}
              -1 &3  &1 \\
               4 &5  &4 \\
               5 &-5 &-1
            \end{vmat}
            =+12
          \end{equation*}
          起始箱体是正定向的，该变换
          保持定向（因为矩阵的行列式为
          正），而终止箱体也是正定向的。
      \end{exparts}
    \end{answer}
  \item 
    每个变换把箱体的大小改变了多少倍？
    \begin{exparts*}
      \partsitem $\colvec{x \\ y}\mapsto\colvec{2x \\ 3y}$
      \partsitem $\colvec{x \\ y}\mapsto\colvec{3x-y \\ -2x+y}$
      \partsitem $\colvec{x \\ y \\ z}\mapsto\colvec{x-y \\ x+y+z \\ y-2z}$
    \end{exparts*}
    \begin{answer}
      将每个变换关于标准基表出，并求出行列式。
      \begin{exparts*}
        \partsitem $6$
        \partsitem $-1$
        \partsitem $-5$
      \end{exparts*}
    \end{answer}
  \item 
    矩形 \( [2..4]\times [2..5] \) 在
    这个矩阵作用下的像的面积是多少？
    \begin{equation*}
       \begin{mat}[r]
         2  &3  \\
         4  &-1
       \end{mat}
    \end{equation*}
    \begin{answer}
      起始面积为 \( 6 \)，该矩阵把大小改变
      \( -14 \) 倍。
      因此像的面积为 \( 84 \)。  
    \end{answer}
  \item
     若 \( \map{t}{\Re^3}{\Re^3} \) 把体积改变 \( 7 \) 倍，
     \( \map{s}{\Re^3}{\Re^3} \) 把体积改变 \( 3/2 \) 倍，
     那么它们的复合把体积改变多少倍？
     \begin{answer}
        改变 \( 21/2 \) 倍。
     \end{answer}
  \item 
    箱体的定义与张成的定义在哪些方面不同？
    \begin{answer}
      对于箱体，我们取的是一个向量的序列（如注中所说，
      向量的顺序是重要的），
      而对于张成，我们取的是一个向量的集合。
      此外，$\Re^n$ 的箱体子集必须有 $n$ 个向量； 
      当然，对于张成，向量的个数可以是任意的。
      最后，对于箱体，系数 $t_1$、~\ldots、$t_n$
      取自区间 $[0..1]$，而对于
      张成，系数可以取遍整个 $\Re$。 
    \end{answer}
  \item 
    \( \deter{TS}=\deter{ST} \) 吗？
    \( \deter{T(SP)}=\deter{(TS)P} \) 吗？
    \begin{answer}
      两问的答案都是肯定的。
      例如，第一问是 \( \deter{TS}=\deter{T}\cdot\deter{S}=
                     \deter{S}\cdot\deter{T}=\deter{ST} \)。  
    \end{answer}
  \item 证明：不存在满足下面条件的 $\nbyn{2}$ 矩阵 $A$ 与~$B$。 
   % http://math.stackexchange.com/questions/827262/need-help-with-a-linear-algebra-proof
   \begin{equation*}
     AB=\begin{mat}[r]
       1 &-1 \\
       2  &0
     \end{mat}
     \quad
     BA=\begin{mat}[r]
       2 &1 \\
       1  &1
     \end{mat}
   \end{equation*}
   \begin{answer}  % due to math.stackexchange.com user dgrasines517
     因为 $\deter{AB}=\deter{A}\cdot\deter{B}=\deter{BA}$，而这两个
     矩阵的行列式不同。
   \end{answer}
  \item 
   \begin{exparts}
     \partsitem 设 \( \deter{A}=3 \) 且 \( \deter{B}=2 \)。
        求 \( \deter{A^2\cdot \trans{B}\cdot B^{-2}\cdot \trans{A} } \)。
     \partsitem 设 \( \deter{A}=0 \)。
        证明 \( \deter{6A^3+5A^2+2A}=0 \)。
    \end{exparts}
    \begin{answer}
      \begin{exparts}
        \partsitem 若它有定义，则它为
           \( (3^2)\cdot (2)\cdot (2^{-2})\cdot (3) \)。
        \partsitem \( \deter{6A^3+5A^2+2A}=\deter{A}\cdot\deter{6A^2+5A+2I} \)。
      \end{exparts}  
    \end{answer}
  \recommended \item
    设 \( T \) 是表示（关于标准基）将平面向量逆时针旋转
    \( \theta \) 弧度的映射的矩阵。
    \( T \) 把大小改变多少倍？
    \begin{answer}
       \(\begin{vmat}
                \cos\theta  &-\sin\theta  \\
                \sin\theta  &\cos\theta
              \end{vmat}=1 \)  
    \end{answer}
  \recommended \item
    保持面积的变换 \( \map{t}{\Re^2}{\Re^2} \) 是否也必须
    保持长度？
    \begin{answer}
      不是，例如
      \begin{equation*}
        T=\begin{mat}[r]
          2  &0  \\
          0  &1/2
        \end{mat}
      \end{equation*}
      的行列式为 \( 1 \)，所以它保持面积，但向量 \( T\vec{e}_1 \)
      的长度为 \( 2 \)。  
    \end{answer}
  \item
    由一个线性相关集所界定的 \( \Re^3 \) 中平行六面体的
    体积是多少？
    \begin{answer}
       为零。  
    \end{answer}
  \recommended \item
    求 \( \Re^3 \) 中以 \( (1,2,1) \)、\( (3,-1,4) \) 与 \( (2,2,2) \)
    为端点的三角形的面积。
    （这里问的是面积，而不是体积。
    该三角形确定了一个平面；问的是这个平面内该三角形的面积。）
    \begin{answer}
      三角形三条边中有两条由这些向量构成。
      \begin{equation*}
        \colvec[r]{2 \\ 2 \\ 2}-\colvec[r]{1 \\ 2 \\ 1}=\colvec[r]{1 \\ 0 \\ 1}
        \qquad
        \colvec[r]{3 \\ -1 \\ 4}-\colvec[r]{1 \\ 2 \\ 1}=\colvec[r]{2 \\ -3 \\ 3}
      \end{equation*}
      求这个三角形面积的一种方法是：先求由这两个向量与一个
      与它们都正交的单位向量所形成的平行六面体的体积，
      再取其一半。
      为求与那两个向量都正交的向量族，可以从两个关系式
      \begin{equation*}
        \colvec[r]{1 \\ 0 \\ 1}
        \dotprod\colvec{x \\ y \\ z}
        =\colvec[r]{0 \\ 0 \\ 0}
        \qquad
        \colvec[r]{2 \\ -3 \\ 3}
        \dotprod\colvec{x \\ y \\ z}
        =\colvec[r]{0 \\ 0 \\ 0}
      \end{equation*}
      出发，解方程组
      \begin{equation*}
        \begin{linsys}{3}
          x  &   &   &+  &z  &=  &0  \\
         2x  &-  &3y &+  &3z &=  &0  
        \end{linsys}
        \grstep{-2\rho_1+\rho_2}
        \begin{linsys}{3}
          x  &   &   &+  &z  &=  &0  \\
             &   &-3y&+  &z  &=  &0  
        \end{linsys}
      \end{equation*}
      得到这个解集。
      \begin{equation*}
        \set{\colvec[r]{-1 \\ 1/3 \\ 1}z\suchthat z\in\Re}
      \end{equation*}
      下面是一个长度为一的解。
      \begin{equation*}
        \frac{1}{\sqrt{19/9}}\cdot\colvec[r]{-1 \\ 1/3 \\ 1}
        =\colvec[r]{-3/\sqrt{19} \\ 1/\sqrt{19} \\ 3/\sqrt{19}}
      \end{equation*}
      于是三角形的面积是这个行列式
      绝对值的一半。
      \begin{equation*}
        \begin{vmat}[r]
             1  &2   &-3/\sqrt{19}   \\
             0  &-3  &1/\sqrt{19}   \\
             1  &3   &3/\sqrt{19}
        \end{vmat}
        =-19/\sqrt{19}
      \end{equation*} 
      绝对值的一半是 $\sqrt{19}/2$。
    \end{answer}
  \item \label{exer:DetProdEqProdDetsFcn}
    \nearbytheorem{th:MatChVolByDetMat} 的另一种证明
    利用了行列式函数的定义。
    \begin{exparts}
      \partsitem 注意：构成
        $S$ 的那些向量构成线性相关集当且仅当
        $\deter{S}=0$，并验证在这种情况下结论成立。
      \partsitem 对于 $\deter{S}\neq 0$ 的情形，为证
        对所有变换都有
        $\deter{TS}/\deter{S}=\deter{T}$，考虑函数
        \( \map{d}{\matspace_{\nbyn{n}}}{\Re} \)，其定义为
        \( T\mapsto \deter{TS}/\deter{S} \)。
        证明 $d$ 具有行列式的第一条性质。
      \partsitem 证明 $d$ 具有行列式函数
        其余的三条性质。
      \partsitem 由此得出 $\deter{TS}=\deter{T}\cdot\deter{S}$。 
    \end{exparts}
    \begin{answer}
      \begin{exparts}
        \partsitem 因为线性相关集的像是
          线性相关的，
          所以若构成 $S$ 的向量构成线性相关集，
          即 $\deter{S}=0$，
          则构成 $t(S)$ 的向量构成线性相关集，
          从而 $\deter{TS}=0$，此时等式成立。
        \partsitem 我们必须验证：若
          $T\smash[b]{\grstep{k\rho_i+\rho_j}}\hat{T}$，则
          $d(T)=\deter{TS}/\deter{S}=\deter{\hat{T}S}/\deter{S}=d(\hat{T})$。
          为此可以验证：先把行做倍加再相乘得到 \( \hat{T}S \)，与
          先相乘得到 \( TS \) 再做倍加，
          结果相同（因为行列式 \( \deter{TS} \) 不受行倍加的
          影响，
          这样便有 \( \deter{\hat{T}S}=\deter{TS} \)，从而
          \( d(\hat{T})=d(T) \)）。
          具体验证如下：~把 \( TS \) 的第~\( i \) 行的
          \( k \) 倍加到
          \( TS \) 的第~$j$ 行之后，其 \( j,p \) 元为
          \( (kt_{i,1}+t_{j,1})s_{1,p}+\dots+(kt_{i,r}+t_{j,r})s_{r,p} \)，
          这正是 \( \hat{T}S \) 的 \( j,p \) 元。
        \partsitem 对于第二条性质，只需验证：对换
          $T\smash[b]{\grstep{\rho_i\swap\rho_j}}\hat{T}$
          后再相乘得到 \( \hat{T}S \)，与先把 \( T \) 乘以 \( S \) 再做对换，
          结果相同（因为行列式 \( \deter{TS} \) 在
          行对换下变号，
          这样便有 \( \deter{\hat{T}S}=-\deter{TS} \)，
          从而 \( d(\hat{T})=-d(T) \)）。 
          该验证与第一条性质的验证完全类似。

          对于第三条性质，只需证明：先做
          $T\smash[b]{\grstep{k\rho_i}}\hat{T}$
          再计算 \( \hat{T}S \)，与先计算 \( TS \) 再做那次
          倍乘，
          结果相同（由于行列式 \( \deter{TS} \) 被缩放 \( k \) 倍，
          我们将有 \( \deter{\hat{T}S}=k\deter{TS} \)，
          从而 \( d(\hat{T})=k\,d(T) \)）。
          这里的论证同样与上面完全类似。

          第四条性质，即若 $T$ 为 $I$ 则结果为 $1$，
          是显然的。
        \partsitem 行列式函数是唯一的，所以
          \( \deter{TS}/\deter{S}=d(T)=\deter{T} \)，
          从而 $\deter{TS}=\deter{T}\deter{S}$。
      \end{exparts}
    \end{answer}
%  \item  
%    Use the fact that
%    \( \deter{TS}=\deter{T}\,\deter{S} \)
%    to prove that if \( \phi \) and \( \sigma \) are \( n \)-permutations
%    then \( \sgn(\phi\sigma)=\sgn(\phi)\sgn(\sigma) \).
%    \cite{HoffmanKunze}
%    \begin{answer}
%       Take \( T=P_\phi \) and \( S=P_\sigma \).
%      Note that \( TS=P_{\phi\sigma} \) and so
%      \( \deter{P_{\phi\sigma}}=\deter{P_\phi}\cdot\deter{P_\sigma} \).  
%    \end{answer}
  \item 
    给出一个满足性质 \( \trans{A}=A^{-1} \) 的非单位矩阵。
    证明：若 \( \trans{A}=A^{-1} \)，则 \( \deter{A}=\pm 1 \)。
    逆命题成立吗？
    \begin{answer}
      任何置换矩阵都具有矩阵的转置等于其逆的性质。

      对于这个蕴含关系，我们知道 \( \deter{\trans{A}}=\deter{A} \)。
      于是 \( 1=\deter{A\cdot A^{-1}}=\deter{A\cdot\trans{A}}
               =\deter{A}\cdot\deter{\trans{A}}=\deter{A}^2 \)。

      逆命题不成立；下面是一个反例。
      \begin{equation*}
        \begin{mat}[r]
          3  &1  \\
          2  &1
        \end{mat}
      \end{equation*}
    \end{answer}
  \item 
    行列式的一条代数
    性质是：从单独一行中提出一个标量因子会使行列式乘上这个标量，这表明
    当 \( H \) 为
    \( \nbyn{3} \) 矩阵时，\( cH \) 的行列式是 \( H \) 的
    行列式的 \( c^3 \) 倍。
    请从几何上解释这一点，也就是
    用 \nearbytheorem{th:MatChVolByDetMat} 来解释。
    （把一个三维物体的线性尺寸放大 $c$ 倍，其体积将放大
    $c^3$ 倍，而其表面积只增加与
    $c^2$ 成比例的量，这一观察就是
    \definend{平方-立方定律}~\cite{Wikipedia}。）
    \begin{answer}
      当箱体的边长变为原来的 \( c \) 倍时，箱体的
      体积就有原来的 \( c^3 \) 倍立方单位。  
    \end{answer}
  \item 
    若存在非奇异矩阵 $P$ 使得 $H=P^{-1}GP$，则称矩阵 $H$ 与 $G$ 是
    \definend{相似的}\index{similar}\index{matrix!similar}
    （我们将在第五章研究这一关系）。
    证明：相似矩阵有相同的行列式。
    \begin{answer}
      若 \( H=P^{-1}GP \)，
      则 \( \deter{H}=\deter{P^{-1}}\deter{G}\deter{P}
        =\deter{P^{-1}}\deter{P}\deter{G}=\deter{P^{-1}P}\deter{G}
        =\deter{G} \)。  
    \end{answer}
  \item  \label{exer:BasisOrient}
    我们通常关于标准基表示 \( \Re^2 \) 中的向量，因此第一象限中的
    向量两个坐标都为正。
    \begin{center}
      \parbox{.75in}{\hbox{}\hfil\includegraphics{det/mp/ch4.45}\hfil\hbox{}}
      \qquad
      \( \rep{\vec{v}}{\stdbasis_2}=\colvec[r]{+3 \\ +2} \)
    \end{center}
    沿逆时针方向绕原点移动，我们依次经过四个区域：
    {\scriptsize
    \begin{equation*}
       \cdots
       \;\longrightarrow\colvec{+ \\ +}
       \;\longrightarrow\colvec{- \\ +}
       \;\longrightarrow\colvec{- \\ -}
       \;\longrightarrow\colvec{+ \\ -}
       \;\longrightarrow\cdots\,.
    \end{equation*} }
    使用这个基
    \begin{center}
      \( B=\sequence{\colvec[r]{0 \\ 1},\colvec[r]{-1 \\ 0}} \)
      \qquad
      \parbox{.75in}{\hbox{}\hfil\includegraphics{det/mp/ch4.46}\hfil\hbox{}}
    \end{center}
    也给出同样的逆时针循环。
    我们说这两个基具有相同的\emph{定向}。\index{orientation}
    \begin{exparts}
      \partsitem 为什么它们给出相同的循环？
      \partsitem 轴上单位向量的其他哪些配置给出
        相同的循环？
      \partsitem 求由那些（有序的）基构成的矩阵的行列式。
      \partsitem 还可能出现哪些逆时针循环，
        相应的
        行列式是什么？
      \partsitem 在 \( \Re^1 \) 中会发生什么？
      \partsitem 在 \( \Re^3 \) 中会发生什么？
    \end{exparts}
    关于定向的一段面向普通读者的精彩讨论见
    \cite{Gardner}。
    \begin{answer}
      \begin{exparts}
        \partsitem 新基是旧基旋转 \( \pi/4 \) 得到的。
        \partsitem 
          $
             \sequence{\colvec[r]{-1 \\ 0},
                       \colvec[r]{0 \\ -1}}
          $、$
             \sequence{\colvec[r]{0 \\ -1},
                       \colvec[r]{1 \\ 0}}
          $
        \partsitem 每种情形下行列式都是 \( +1 \)
          （我们说这些基
          具有\definend{正定向}）。
        \partsitem 因为一次只能改变一个符号，仅有的另一种
          可能的循环是
          \begin{equation*}
             \cdots
             \;\longrightarrow\colvec{+ \\ +}
             \;\longrightarrow\colvec{+ \\ -}
             \;\longrightarrow\colvec{- \\ -}
             \;\longrightarrow\colvec{- \\ +}
             \;\longrightarrow\cdots\,.
          \end{equation*}
          这里每个相应的行列式都是 \( -1 \)
          （我们说这样的基具有\definend{负定向}）。
        \partsitem 有一个正定向的基 \( \sequence{(1)} \)
          和一个负定向的基 \( \sequence{(-1)} \)。
        \partsitem 共有 \( 48 \) 个基（第一个单位向量有
          \( 6 \) 种半轴可选，
          第二个有 \( 4 \) 种，最后一个有 \( 2 \) 种）。
          其中一半像下面左图中的标准基那样是正定向的，
          另一半像右图中的那样是负定向的
         \begin{center}
           \includegraphics{det/mp/ch4.47}
           \hspace*{4em}
           \includegraphics{det/mp/ch4.48}
          \end{center}
          在 \( \Re^3 \) 中，正定向有时被称为
          “右手定向”，因为如果一个人伸出
          右手，
          手指从 \( \vec{e}_1 \) 卷向 \( \vec{e}_2 \)，
          那么拇指就会指向 \( \vec{e}_3 \)。
      \end{exparts}  
    \end{answer}
%  \item 
%    A region of \( \Re^n \) is \definend{convex}\index{convex region} 
%    if for any two points
%    connecting that region, the line segment joining them lies entirely
%    inside the region (the inside of a sphere is convex, while the skin of a
%    sphere or a horseshoe is not).
%    Prove that boxes are convex.
%    \begin{answer}
%      Let \( \vec{p}=p_1\vec{v}_1+\dots +p_n\vec{v}_n \) and
%      \( \vec{q}=q_1\vec{v}_1+\dots +q_n\vec{v}_n \) be two vectors from a box
%      so that \( p_1,\dots,\,p_n,q_1,\dots,\,q_n\in [0..1] \).
%      The line segment between them is this.
%      \begin{equation*}
%        \set{t\cdot\vec{p}+(1-t)\cdot\vec{q}
%              \suchthat t\in [0..1]}
%        =\set{(tp_1+(1-t)q_1)\cdot\vec{v}_1+\dots+(tp_n+(1-t)q_n)\cdot\vec{v}_n
%              \suchthat t\in [0..1] }
%      \end{equation*}
%      Showing that each member of that set is in the box is routine.  
%    \end{answer}
  \item \label{exer:DetProdEqProdDetsPerms}
    \textit{本题用到选读小节“行列式函数的存在性”中的内容。}
    利用行列式的置换展开公式证明
    \nearbytheorem{th:MatChVolByDetMat}。
    \begin{answer}
      我们将比较 \( \det(\vec{s}_1,\dots,\vec{s}_n) \) 与
      \( \det(t(\vec{s}_1),\dots,t(\vec{s}_n)) \)，以证明后者
      与前者相差一个因子 $\deter{T}$。
      我们关于标准基表示那些 \( \vec{s}\, \)：
      \begin{equation*}
        \rep{\vec{s}_i}{\stdbasis_n}=
           \colvec{s_{1,i} \\ s_{2,i} \\ \vdots \\ s_{n,i}}
      \end{equation*}
      然后用矩阵—向量乘法
      表示映射的作用
      \begin{align*}
        \rep{\,t(\vec{s}_i)\,}{\stdbasis_n}
         &=\generalmatrix{t}{n}{n}
           \colvec{s_{1,j} \\ s_{2,j} \\ \vdots \\ s_{n,j}}     \\
         &=s_{1,j}\colvec{t_{1,1} \\ t_{2,1} \\ \vdots \\ t_{n,1}}
          +s_{2,j}\colvec{t_{1,2} \\ t_{2,2} \\ \vdots \\ t_{n,2}}
          +\dots
          +s_{n,j}\colvec{t_{1,n} \\ t_{2,n} \\ \vdots \\ t_{n,n}} \\
         &=s_{1,j}\vec{t}_1+s_{2,j}\vec{t}_2+\dots+s_{n,j}\vec{t}_n
      \end{align*}
      其中 \( \vec{t}_i \) 是 \( T \) 的第~$i$ 列。
      于是 $\det(t(\vec{s}_1),\,\dots,\,t(\vec{s}_n))$ 等于
      $
        \det(s_{1,1}\vec{t}_1\!+\!s_{2,1}\vec{t}_2\!
                +\!\dots\!+\!s_{n,1}\vec{t}_n,\,
             \dots,\,
             s_{1,n}\vec{t}_1\!+\!s_{2,n}\vec{t}_2
                \!+\!\dots\!+\!s_{n,n}\vec{t}_n)
     $。

     如同置换展开公式的推导那样，我们
     利用多重线性，
     先沿第一个变元中的和式拆分
     \begin{multline*}
         \det(s_{1,1}\vec{t}_1,\,
            \dots,\,
            s_{1,n}\vec{t}_1+s_{2,n}\vec{t}_2+\dots+s_{n,n}\vec{t}_n)  \\ 
        +\cdots{}                                             
        +\det(s_{n,1}\vec{t}_n,\,
           \ldots,\,
            s_{1,n}\vec{t}_1+s_{2,n}\vec{t}_2+\dots+s_{n,n}\vec{t}_n)
     \end{multline*}
     再沿第二个变元中的和式拆分这 $n$ 个和式中的每一个，如此等等。
     与置换展开的推导一样，我们最终得到
     \( n^n \) 个和式行列式，每个形如
     $\det(s_{i_1,1}\vec{t}_{i_1},s_{i_2,2}\vec{t}_{i_2},
            \,\dots,\,
            s_{i_n,n}\vec{t}_{i_n})$。
     把每个 $s_{i,j}$ 提出来，得
     $=
       s_{i_1,1}s_{i_2,2}\dots s_{i_n,n}
        \cdot\det(\vec{t}_{i_1},\vec{t}_{i_2},
       \,\dots,\,
       \vec{t}_{i_n})
      $.

      如同置换展开的推导那样，
      只要 $i_1$、\ldots、$i_n$ 中有两个指标相等，
      该行列式
      就有两个相同的变元，其值为 $0$。
      因此只需考虑 $i_1$、\ldots、$i_n$ 构成数
      $1$、\ldots、$n$ 的一个置换的情形。
      于是我们有
      \begin{equation*}
        \det(t(\vec{s}_1),\dots,t(\vec{s}_n))=
          \sum_{\text{permutations\ } \phi}
            s_{\phi(1),1}\dots s_{\phi(n),n}
            \det(\vec{t}_{\phi(1)},\dots,\vec{t}_{\phi(n)}).
      \end{equation*}
      把 $\det(\vec{t}_{\phi(1)},\ldots,\vec{t}_{\phi(n)})$
      中的列对换，便还原出矩阵 \( T \)，这会按因子
      $\sgn{\phi}$ 改变符号，
      然后再提出 $T$ 的行列式。
      \begin{equation*}
        =\sum_\phi
          s_{\phi(1),1}\dots s_{\phi(n),n}
            \det(\vec{t}_1,\dots,\vec{t}_n)\cdot\sgn{\phi}
        =\det(T)\sum_\phi
          s_{\phi(1),1}\dots s_{\phi(n),n}\cdot\sgn{\phi}.
      \end{equation*}
      如同在证明矩阵的行列式
      等于其转置的
      行列式时那样，我们交换各个 $s$ 的位置，改为按行号递增
      （而不是按列号递增）的顺序列出它们
      （并把 $\sgn(\phi^{-1})$ 代入 $\sgn(\phi)$ 的位置）。
      \begin{equation*}
        =\det(T)\sum_\phi
          s_{1,\phi^{-1}(1)}\dots s_{n,\phi^{-1}(n)}\cdot\sgn{\phi^{-1}}  
        =\det(T)\det(\vec{s}_1,\vec{s}_2,\dots,\vec{s}_n)
      \end{equation*}
    \end{answer}
%  \item 
%    Suppose that \( \map{f}{\matspace_{\nbyn{n}}}{\Re} \) is a non-constant
%    function with the property that \( f(GH)=f(G)f(H) \).
%    \begin{exparts}
%      \partsitem Show that \( f \) sends the identity to \( 1 \).
%      \partsitem Show that \( f \) maps the elementary matrix 
%        \( C_{i,j}(k) \) to \( 1 \)
%        (this matrix results from performing \( k\rho_i+\rho_j \) to the
%        identity).
%      \partsitem Show that 
%        \( f \) maps a row swap matrix to \( +1 \) or \( -1 \).
%    \end{exparts}
%    \begin{answer}
%      \begin{exparts}
%        \partsitem We have \( f(IH)=f(I)f(H) \) and \( f(IH)=f(H) \).
%        \partsitem
%        \partsitem A row swap matrix has the property that when 
%          done twice it equals
%          the identity.
%          But \( f(R)f(R)=f(RR)=f(I)=1 \) implies that \( f(R)=\pm 1 \).
%      \end{exparts} 
%     \end{answer}
  \recommended \item
    \begin{exparts}
      \partsitem 证明下式给出 \( \Re^2 \) 中过点
        \( (x_2,y_2) \) 与 \( (x_3,y_3) \) 的
        直线方程。
        \begin{equation*}
          \begin{vmat}
            x    &x_2 &x_3  \\
            y    &y_2 &y_3  \\
            1    &1   &1
          \end{vmat}=0
        \end{equation*}
      \partsitem \cite{Monthly55p249}
        证明：以 \( (x_1,y_1) \)、
        \( (x_2,y_2) \) 与 \( (x_3,y_3) \) 为顶点的三角形的面积为
        \begin{equation*}
          \frac{1}{2}
          \begin{vmat}
            x_1  &x_2 &x_3  \\
            y_1  &y_2 &y_3  \\
            1    &1   &1
          \end{vmat}.
        \end{equation*}
      \partsitem \cite{MathMag73p286}
        证明：顶点位于 \( (x_1,y_1) \)、
        \( (x_2,y_2) \) 与 \( (x_3,y_3) \)，且坐标均为整数的三角形，
        其面积必为某个正整数 \( N \) 的 \( N \) 或 \( N/2 \)。
    \end{exparts}
    \begin{answer}
      \begin{exparts}
        \partsitem 代数上的验证是容易的。
        \begin{equation*}
          0
          =xy_2+x_2y_3+x_3y-x_3y_2-xy_3-x_2y 
          =x\cdot (y_2-y_3)+y\cdot (x_3-x_2)+x_2y_3-x_3y_2 
        \end{equation*}
        化简后得到熟悉的形式
        \begin{equation*}
          y=x\cdot (x_3-x_2)/(y_3-y_2)+(x_2y_3-x_3y_2)/(y_3-y_2)
        \end{equation*}
        （$y_3-y_2=0$ 的情形也容易处理）。

        从几何上看，这幅
        图显示了由这三个向量所形成的箱体。
        注意，所有
        三个向量的终点都在 $z=1$ 平面内。
        右边两个向量的下方是过点
        $(x_2,y_2)$ 与 $(x_3,y_3)$ 的直线。
        \begin{center}
          \includegraphics{det/mp/ch4.49}
        \end{center}
        该箱体
        将有非零的体积，除非这三个向量的终点所构成的三角形
        是退化的。
        这只会在
        （假定 $(x_2,y_3)\neq (x_3,y_3)$ 时）
        $(x,y)$ 位于过另两点的直线上时发生。 
       \partsitem \answerasgiven %
        我们按通常的办法，从上式的法向形式出发，
        求以 $(x_1,y_1)$、$(x_2,y_2)$ 与 $(x_3,y_3)$ 为顶点的
        三角形过 $(x_1,y_1)$ 的高：
        \begin{equation*}
          \frac{1}{\sqrt{(x_2-x_3)^2+(y_2-y_3)^2}}
          \begin{vmat}
            x_1  &x_2  &x_3  \\
            y_1  &y_2  &y_3  \\
            1    &1    &1
          \end{vmat}.
        \end{equation*}
        再走一步就得到三角形的面积为
        \begin{equation*}
          \frac{1}{2}
          \begin{vmat}
            x_1  &x_2  &x_3  \\
            y_1  &y_2  &y_3  \\
            1    &1    &1
          \end{vmat}.
        \end{equation*}
        与通常那种证明一列项与行列式恒等的办法相比，
        这里的阐述更清楚地揭示了\textit{modus operandi}（做法）。
       \partsitem  \answerasgiven %
        设
        \begin{equation*}
          D=
          \begin{vmat}
            x_1  &x_2  &x_3  \\
            y_1  &y_2  &y_3  \\
            1    &1    &1
          \end{vmat}
        \end{equation*}
        则三角形的面积为 $(1/2)\deter{D}$。
        于是，若坐标全为整数，则 $D$ 是一个整数。
      \end{exparts}
    \end{answer}
\end{exercises}
